\documentclass[12pt]{article}
\usepackage[utf8]{inputenc}
\usepackage[T1]{fontenc}
\usepackage[margin=1in]{geometry}
\usepackage{setspace}
\setlength{\parindent}{2em}
\setlength{\parskip}{1.25em}
\renewcommand{\baselinestretch}{1.0}
\makeatletter
\newcommand{\@subtitle}{}
\newcommand{\subtitle}[1]{\renewcommand{\@subtitle}{#1}}
\renewcommand{\@maketitle}{%
  \newpage
  \null
  \vskip 1em%
  \begin{center}%
    {\LARGE \@title \par}%
    \ifx\@subtitle\@empty\else
      \vskip 0.5em%
      {\large \@subtitle \par}%
    \fi
  \end{center}%
  \par
  \vskip 1em}
\makeatother
\usepackage{xcolor}
\usepackage[colorlinks=true,urlcolor=blue!55!black,linkcolor=black]{hyperref}

\title{\textit{Allen v.\ Wright}, 468 U.S.\ 737 (1984)}
\subtitle{Case Brief}
\author{}
\date{}

\begin{document}
\maketitle

\section{Memory Jogger}
Parents of Black public-school children lacked Article~III standing to sue the IRS over its lax enforcement of rules denying tax-exempt status to racially discriminatory private schools because abstract stigmatic harm without personal denial of equal treatment is not cognizable and any impairment of public-school desegregation is not ``fairly traceable'' to the IRS, grounded in separation-of-powers limits on judicial oversight of the Executive (\href{https://supreme.justia.com/cases/federal/us/468/737/}{Justia}; \href{https://www.oyez.org/cases/1983/81-757}{Oyez}).

\section{Detailed Case Facts}
Following \textit{Brown v.\ Board of Education} (1954) and federal court desegregation orders in the 1960s and 1970s, numerous private ``segregation academies'' opened or expanded in communities undergoing public-school desegregation. Prompted by litigation in Mississippi (\textit{Green v.\ Connally}, 330 F.\ Supp.\ 1150 (D.D.C.\ 1971), summarily aff'd sub nom.\ \textit{Coit v.\ Green}, 404 U.S.\ 997 (1971)), the Internal Revenue Service (IRS) adopted a policy in 1970--1971 (Rev.\ Rul.\ 71-447) denying tax-exempt status under Internal Revenue Code \S\S~501(a) and 501(c)(3)---and charitable contribution deductibility under \S\S~170(a)(1) and 170(c)(2)---to racially discriminatory private schools, a statutory interpretation later upheld in \href{https://supreme.justia.com/cases/federal/us/461/574/}{\textit{Bob Jones University v.\ United States}, 461 U.S.\ 574 (1983)}. To enforce that policy, the IRS issued Rev.\ Proc.\ 75-50 (1975), requiring schools to adopt, publicize, and annually certify a racially nondiscriminatory policy (\href{https://supreme.justia.com/cases/federal/us/468/737/}{468 U.S.\ at 740--743}).

In 1976, Inez Wright and other parents of Black children attending public schools in seven States in districts undergoing desegregation filed a nationwide class action on behalf of several million parents and children against the Secretary of the Treasury (Donald Regan) and the Commissioner of Internal Revenue. W.\ Wayne Allen, head of one of the private school systems named in the complaint (Briarcrest Christian Academy in Memphis), intervened as a defendant (\href{https://supreme.justia.com/cases/federal/us/468/737/}{468 U.S.\ at 743--747}; \href{https://www.oyez.org/cases/1983/81-757}{Oyez}). Plaintiffs alleged that the IRS's 1975 guidelines were so ineffective that they allowed racially discriminatory private schools (including 17 named schools or school systems) to obtain tax exemptions merely by certifying a paper policy without actually implementing it. Crucially, plaintiffs never alleged that their children had applied or would ever apply to any private school. Instead, they alleged two injuries: (1) direct harm and racial stigma from the mere fact of federal financial aid to discriminatory private schools; and (2) impairment of their children's opportunity to receive a desegregated public education because tax exemptions subsidized white flight to segregated private schools. They sought declaratory relief and a sweeping injunction forcing the IRS to replace Rev.\ Proc.\ 75-50 with stricter presumptions against tax exemptions (\href{https://supreme.justia.com/cases/federal/us/468/737/}{468 U.S.\ at 745--747}).

\textbf{Key Constitutional and Statutory Provisions:}
\begin{itemize}
  \item \textbf{U.S.\ Const.\ art.\ III, \S~2, cl.\ 1:} Confines the federal judicial power to actual ``Cases'' and ``Controversies,'' from which the constitutional standing requirements (injury in fact, traceability, and redressability) derive.
  \item \textbf{U.S.\ Const.\ art.\ II, \S~3 (Take Care Clause):} Assigns to the Executive Branch the duty to ``take Care that the Laws be faithfully executed.''
  \item \textbf{Internal Revenue Code, 26 U.S.C.\ \S\S~501(a), 501(c)(3), 170(a)(1), 170(c)(2):} Govern tax-exempt charitable status and tax-deductible charitable contributions. Plaintiffs also invoked Title~VI of the Civil Rights Act of 1964 (42 U.S.C.\ \S~2000d et seq.), 42 U.S.C.\ \S~1981, and the Fifth and Fourteenth Amendments.
\end{itemize}

\textbf{Historical Context:} While the lawsuit was pending in 1978--1979, the IRS itself proposed tougher Revenue Procedures to screen out segregation academies, but Congress blocked implementation through appropriations riders (the Ashbrook and Dornan Amendments) for fiscal years 1980--1982 (\href{https://supreme.justia.com/cases/federal/us/468/737/}{468 U.S.\ at 747--748 \& n.16}). The case reached the Supreme Court amidst intense Reagan-era debate over federal judicial oversight of executive agencies and civil-rights enforcement.

\section{Procedural History}
\begin{itemize}
  \item \textbf{U.S.\ District Court for the District of Columbia:} After Congress enacted the 1979 appropriations riders, the district court granted defendants' motion to dismiss (\textit{Wright v.\ Miller}, 480 F.\ Supp.\ 790 (D.D.C.\ 1979)), holding that plaintiffs lacked standing, that the requested injunction would be unduly intrusive, and that judicial relief would contravene Congress's 1979 spending restrictions (\href{https://supreme.justia.com/cases/federal/us/468/737/}{468 U.S.\ at 748}).
  \item \textbf{U.S.\ Court of Appeals for the D.C.\ Circuit:} A divided panel reversed (\textit{Wright v.\ Regan}, 656 F.2d 820 (D.C.\ Cir.\ 1981)), holding (over Judge Tamm's dissent) that plaintiffs had standing based on the ``denigration'' injury suffered by Black parents and schoolchildren when the government grants tax-exempt status to discriminatory schools in their communities, distinguishing \href{https://supreme.justia.com/cases/federal/us/426/26/}{\textit{Simon v.\ Eastern Kentucky Welfare Rights Org.}, 426 U.S.\ 26 (1976)}, and enjoining the IRS pending remand.
  \item \textbf{U.S.\ Supreme Court:} Both the Government (No.\ 81-970, \textit{Regan v.\ Wright}) and intervenor Allen (No.\ 81-757, \textit{Allen v.\ Wright}) petitioned for certiorari. The Supreme Court granted certiorari (462 U.S.\ 1130 (1983)), heard oral argument on February~29, 1984, and on July~3, 1984 \textbf{reversed} the D.C.\ Circuit and vacated its injunction (\href{https://supreme.justia.com/cases/federal/us/468/737/}{468 U.S.\ at 750, 766}).
\end{itemize}

\section{Judicial Votes}
\textbf{Judgment:} 5--3 reversing the D.C.\ Circuit and vacating its injunction (\href{https://supreme.justia.com/cases/federal/us/468/737/}{468 U.S.\ at 739, 766}).
\begin{itemize}
  \item \textbf{Majority Opinion:} Justice O'Connor delivered the opinion of the Court, joined by Chief Justice Burger and Justices White, Powell, and Rehnquist (5 Justices).
  \item \textbf{Concurrences:} None.
  \item \textbf{Dissents:} Justice Brennan filed a dissenting opinion; Justice Stevens filed a dissenting opinion, joined by Justice Blackmun (3 Justices dissenting).
  \item \textbf{Recused:} Justice Marshall took no part in the decision of the cases.
\end{itemize}

\section{Holding}
\textbf{Judgment:} Reversed; the D.C.\ Circuit's injunction is vacated and the complaint must be dismissed for lack of Article~III standing.

\textbf{Rule Announced:} Article~III standing---built on the core idea of separation of powers---requires a plaintiff to allege (1) a personal, judicially cognizable injury, (2) fairly traceable to the defendant's allegedly unlawful conduct, and (3) likely to be redressed by the requested relief. Neither an abstract right to have the Government obey the law nor a stigmatic injury from racial discrimination confers standing unless the plaintiff is \textit{personally denied equal treatment}. Furthermore, where a cognizable injury (such as diminished opportunity for a desegregated public education) depends on the independent choices of third parties not before the court (private school administrators and white parents), the causal chain is too attenuated and speculative to satisfy traceability---especially when the requested relief would enlist federal courts to restructure Executive Branch enforcement programs assigned to the President under the Take Care Clause (\href{https://supreme.justia.com/cases/federal/us/468/737/}{468 U.S.\ at 750--761}).

\textbf{Holding vs.\ Dicta:}
\begin{itemize}
  \item \textbf{Holding:} Plaintiffs lack Article~III standing because (a) their direct-harm/stigmatic claim is not a judicially cognizable injury without personal denial of equal treatment, and (b) their diminished-desegregation injury is not fairly traceable to the challenged IRS tax-exemption practices.
  \item \textbf{Dicta:} (1) Footnote~19's conceptual separation of ``fairly traceable'' (link between unlawful conduct and injury) from ``redressability'' (link between injury and requested relief), which later became canonical in \href{https://supreme.justia.com/cases/federal/us/504/555/}{\textit{Lujan v.\ Defenders of Wildlife}, 504 U.S.\ 555 (1992)} (\href{https://constitution.congress.gov/browse/essay/artIII-S2-C1-6-4-1/ALDE_00012995/}{Constitution Annotated}); (2) the Court's summary of prudential standing limits (third-party rights, generalized grievances, zone of interests) at 468 U.S.\ at 751; and (3) the Court's observation that it need not decide whether standing was properly found in \textit{Coit v.\ Green}.
\end{itemize}

\section{Analysis of Opinions}

\subsection{Majority Opinion (Justice O'Connor)}
Justice O'Connor anchored Article~III standing squarely in structural constitutionalism: ``More important, the law of Art.\ III standing is built on a single basic idea --- the idea of separation of powers'' (\href{https://supreme.justia.com/cases/federal/us/468/737/}{468 U.S.\ at 752}). After reviewing constitutional and prudential standing principles, she evaluated each of plaintiffs' two alleged injuries in turn.

First, she rejected plaintiffs' direct-harm and stigmatic-injury claim as non-cognizable. Standing cannot rest on an asserted right to have the Government obey the law (\textit{Schlesinger v.\ Reservists Committee}, \textit{Valley Forge Christian College}). While racial stigma is a grave noneconomic harm, precedent---including \href{https://supreme.justia.com/cases/federal/us/465/728/}{\textit{Heckler v.\ Mathews}, 465 U.S.\ 728 (1984)}, \textit{Moose Lodge No.\ 107 v.\ Irvis}, \textit{O'Shea v.\ Littleton}, and \textit{Rizzo v.\ Goode}---limits standing to persons ``personally denied equal treatment.'' Otherwise, ``[a] black person in Hawaii could challenge the grant of a tax exemption to a racially discriminatory school in Maine,'' turning federal courts into forums for ``concerned bystanders'' (\href{https://supreme.justia.com/cases/federal/us/468/737/}{468 U.S.\ at 755--756}).

Second, Justice O'Connor agreed that plaintiffs' children's diminished ability to receive a desegregated education was ``beyond any doubt'' a judicially cognizable and profoundly serious injury. Nonetheless, it failed the ``fairly traceable'' requirement. Relying heavily on \href{https://supreme.justia.com/cases/federal/us/426/26/}{\textit{Simon v.\ Eastern Kentucky Welfare Rights Org.}}, she found the causal chain fatally attenuated because it depended on multiple independent third-party decisions: it was uncertain how many discriminatory private schools in plaintiffs' communities held exemptions, speculative whether losing tax-exempt status would cause any school to change its admissions policies, and speculative whether white parents would transfer children back to public schools in numbers sufficient to affect integration (\href{https://supreme.justia.com/cases/federal/us/468/737/}{468 U.S.\ at 757--759}). Separation-of-powers concerns reinforced this conclusion: recognizing standing to restructure an agency's enforcement apparatus would make courts ``continuing monitors'' of the Executive, usurping the President's Article~II Take Care duty. Finally, she distinguished \textit{Gilmore v.\ City of Montgomery} (direct equal-access right to public parks), \textit{Norwood v.\ Harrison} (direct violation of a state desegregation decree), and \textit{Coit v.\ Green} (a non-adversarial summary affirmance limited to Mississippi's pre-1970 IRS policy) (\href{https://supreme.justia.com/cases/federal/us/468/737/}{468 U.S.\ at 761--766}).

\subsection{Dissenting Opinion (Justice Brennan)}
Justice Brennan accused the majority of using ``standing to slam the courthouse door'' by invoking ``separation of powers'' as an empty incantation to obscure the concrete civil-rights injury alleged (\href{https://supreme.justia.com/cases/federal/us/468/737/}{468 U.S.\ at 766--767}). Focusing on plaintiffs' second injury---deprivation of an integrated public education---he emphasized that on a Rule~12(b)(6) motion to dismiss, courts must accept the complaint's allegations as true. Contrary to the majority's reading, the complaint specifically named at least 32 discriminatory private schools (including 18 in Memphis, Tennessee alone) receiving tax benefits in desegregating districts (\href{https://supreme.justia.com/cases/federal/us/468/737/}{468 U.S.\ at 775}).

On causation, Justice Brennan argued that \textit{Norwood v.\ Harrison} and \textit{Coit v.\ Green} squarely established that tangible government financial aid to segregated private schools directly impedes public-school desegregation even without proof that a specific child would switch schools. He distinguished \href{https://supreme.justia.com/cases/federal/us/426/26/}{\textit{Simon}} because the plaintiffs there sought services from private hospitals, whereas here plaintiffs sought desegregated public schools directly undermined by federal subsidies. Echoing scholarly critiques, he warned that the Court's causation test was ``no more than a poor disguise for the Court's view of the merits'' (\href{https://supreme.justia.com/cases/federal/us/468/737/}{468 U.S.\ at 782}).

\subsection{Dissenting Opinion (Justice Stevens, joined by Justice Blackmun)}
Justice Stevens distilled the case into three straightforward propositions: injury in fact was undisputed; causation followed from ``elementary economics''; and separation of powers did not bar jurisdiction (\href{https://supreme.justia.com/cases/federal/us/468/737/}{468 U.S.\ at 783--784}). Citing \textit{Regan v.\ Taxation With Representation} and \href{https://supreme.justia.com/cases/federal/us/461/574/}{\textit{Bob Jones University}}, he observed that tax exemptions and deductions function as ``cash grants.'' By basic price theory---``when something becomes more expensive, less of it will be purchased''---revoking tax-exempt status raises the cost of operating segregated private schools, forcing them either to raise tuition (reducing demand) or to integrate, either of which advances public-school desegregation (\href{https://supreme.justia.com/cases/federal/us/468/737/}{468 U.S.\ at 785--788}). Unlike \href{https://supreme.justia.com/cases/federal/us/426/26/}{\textit{Simon}}, where hospitals faced an offsetting ``financial drain'' to provide free care, the economic incentives here ran solely in one direction.

Turning to separation of powers, Justice Stevens faulted the majority for conflating standing (the plaintiff's personal stake) with justiciability and merits questions (whether the law constrains executive discretion). Under \textit{Marbury v.\ Madison}, deciding whether the Constitution or Internal Revenue Code limits IRS enforcement discretion is a classic judicial function that should be resolved on the merits rather than dismissed on an ``untenable'' economic assumption (\href{https://supreme.justia.com/cases/federal/us/468/737/}{468 U.S.\ at 789--795}).

\section{Examples: Future Applications}
\begin{enumerate}
  \item \textbf{Same Side (No Standing):} Two States sue the Department of Homeland Security seeking an injunction requiring federal officers to arrest and prosecute more noncitizens under federal immigration statutes. \textit{Why:} Directly paralleling \textit{Allen}'s separation-of-powers rationale, the Supreme Court held in \href{https://supreme.justia.com/cases/federal/us/599/22-58/}{\textit{United States v.\ Texas}, 599 U.S.\ 670 (2023)}, that plaintiffs lack Article~III standing to challenge Executive Branch arrest and enforcement priorities where they are not themselves the object of enforcement.
  \item \textbf{Same Side (No Standing):} Environmentalists sue the Secretary of the Interior challenging an agency rule limiting Endangered Species Act consultation abroad, alleging that U.S.\ funding of foreign infrastructure projects will harm endangered species they hope to visit ``some day.'' \textit{Why:} Extending \textit{Allen}'s traceability and separation-of-powers framework, \href{https://supreme.justia.com/cases/federal/us/504/555/}{\textit{Lujan v.\ Defenders of Wildlife}, 504 U.S.\ 555 (1992)}, held that speculative third-party actions and generalized grievances about executive law-execution cannot establish Article~III standing.
  \item \textbf{Opposite Side (Standing Exists):} A Black student applies for admission to a tax-exempt private academy, is personally rejected because of her race, and sues challenging the government's discriminatory funding program. \textit{Why:} Unlike the \textit{Allen} plaintiffs, who never applied to any private school, this plaintiff was \textit{personally denied equal treatment}, satisfying both \textit{Allen}'s stigmatic-injury rule and \href{https://supreme.justia.com/cases/federal/us/465/728/}{\textit{Heckler v.\ Mathews}}.
  \item \textbf{Opposite Side (Standing Exists):} Non-minority construction contractors challenge a city ordinance setting aside 10\% of municipal contracts for minority-owned businesses without proving they would have won a specific bid absent the set-aside. \textit{Why:} In \href{https://supreme.justia.com/cases/federal/us/508/656/}{\textit{Northeastern Fla.\ Chapter, Associated General Contractors v.\ City of Jacksonville}, 508 U.S.\ 656 (1993)}, the Court held that being personally subjected to an unequal governmental barrier in competing for a benefit is itself a concrete, directly traceable equal-protection injury.
  \item \textbf{Fence-Sitter:} Liquid-fuel producers sue the EPA challenging its approval of California's electric-vehicle mandate, arguing that forcing automakers to sell more EVs foreseeably reduces gasoline demand and lowers plaintiffs' revenues. \textit{Why:} Although the injury flows through third-party automakers and consumers (triggering \textit{Allen}'s third-party causation skepticism), the Supreme Court held in \textit{Diamond Alternative Energy, LLC v.\ EPA} (2025) (\href{https://supreme.justia.com/cases/federal/us/606/24-7/}{Justia}) that commonsense economic predictability can satisfy traceability and redressability when regulation is explicitly designed to depress demand for the plaintiff's product---illustrating the persistent tension between \textit{Allen}'s skepticism and Justice Stevens's economic-incentive logic.
\end{enumerate}

\section{Critique}
\textbf{Scholarly Criticism and Logical Weaknesses:} \textit{Allen v.\ Wright} is one of the most heavily critiqued standing decisions in modern constitutional law. Scholars such as Gene R.\ Nichol, Jr.\ (\textit{Abusing Standing: A Comment on Allen v.\ Wright}, 133 U.\ Pa.\ L.\ Rev.\ 635 (1985)) and commentary in the \href{https://scholarship.law.unc.edu/cgi/viewcontent.cgi?article=2983&context=nclr}{North Carolina Law Review (1985)} argue that the Court manipulated causation to duck a politically sensitive merits dispute, effectively imposing a summary-judgment burden of proof at the Rule~12(b)(6) pleading stage. Logically, commentators highlight a sharp asymmetry in the Court's economic assumptions: in \textit{Allen}, the Court dismissed as ``pure speculation'' the idea that revoking tax exemptions would alter private-school pricing and white-flight enrollment, yet in subsequent cases (such as \href{https://supreme.justia.com/cases/federal/us/508/656/}{\textit{Northeastern Florida}} and \href{https://supreme.justia.com/cases/federal/us/606/24-7/}{\textit{Diamond Alternative Energy}}) the Court readily credited commonsense market incentives or reframed the injury as loss of an equal opportunity rather than guaranteed outcome.

\textbf{Progressive / Pragmatic Perspective:} Progressive scholars---most prominently \href{https://chicagounbound.uchicago.edu/context/journal_articles/article/12523/viewcontent/1458.pdf}{Cass R.\ Sunstein} and \href{https://scholarship.law.duke.edu/cgi/viewcontent.cgi?article=3221&context=dlj}{Duke Law Journal} commentators---argue that \textit{Allen} institutionalized a deregulatory asymmetry in public-law litigation: regulated entities (such as Bob Jones University) always have standing to challenge aggressive agency enforcement, while regulatory beneficiaries (such as Black public-school families) are routinely locked out when agencies underenforce statutory or constitutional mandates. By folding executive-discretion concerns into Article~III traceability, \textit{Allen} insulated executive non-enforcement of civil-rights laws from judicial review.

\textbf{Originalist / Textualist Perspective:} Conservative formalists like then-Judge Antonin Scalia celebrated \textit{Allen}'s central thesis---that Article~III standing is an essential structural safeguard of separation of powers preventing unelected judges from usurping the President's Article~II Take Care authority and Congress's oversight role (\href{https://lawreview.uchicago.edu/print-archive/justice-scalias-other-standing-legacy}{U.\ Chi.\ L.\ Rev.}). Yet a growing body of originalist scholarship (see \href{https://lawreview.gmu.edu/print__issues/the-court-continues-to-confuse-standing-the-pitfalls-of-faux-article-iii-originalism/}{George Mason Law Review}) critiques the modern ``injury-in-fact / traceability'' trilogy crystallized by \textit{Allen} and \href{https://supreme.justia.com/cases/federal/us/504/555/}{\textit{Lujan}} as ``faux originalism.'' At the Founding, ``Cases'' and ``Controversies'' turned on whether the plaintiff possessed a legally cognizable private right or a valid statutory cause of action (including \textit{qui tam} and relator suits), not on a freestanding, malleable judicial inquiry into sociological causation chains.

\section{Key Quotations}
\begin{quote}
``More important, the law of Art.\ III standing is built on a single basic idea --- the idea of separation of powers.'' --- Justice O'Connor, Opinion of the Court, 468 U.S.\ at 752 (\href{https://supreme.justia.com/cases/federal/us/468/737/}{Justia}).
\end{quote}

\begin{quote}
``If the abstract stigmatic injury were cognizable, standing would extend nationwide to all members of the particular racial groups against which the Government was alleged to be discriminating by its grant of a tax exemption to a racially discriminatory school, regardless of the location of that school. All such persons could claim the same sort of abstract stigmatic injury respondents assert in their first claim of injury. A black person in Hawaii could challenge the grant of a tax exemption to a racially discriminatory school in Maine.'' --- Justice O'Connor, Opinion of the Court, 468 U.S.\ at 755--756 (\href{https://supreme.justia.com/cases/federal/us/468/737/}{Justia}).
\end{quote}

\begin{quote}
``The Constitution, after all, assigns to the Executive Branch, and not to the Judicial Branch, the duty to `take Care that the Laws be faithfully executed.'\,'' --- Justice O'Connor, Opinion of the Court, 468 U.S.\ at 761 (\href{https://supreme.justia.com/cases/federal/us/468/737/}{Justia}).
\end{quote}

\begin{quote}
``More than one commentator has noted that the causation component of the Court's standing inquiry is no more than a poor disguise for the Court's view of the merits of the underlying claims.'' --- Justice Brennan, dissenting, 468 U.S.\ at 782 (\href{https://supreme.justia.com/cases/federal/us/468/737/}{Justia}).
\end{quote}

\begin{quote}
``This causation analysis is nothing more than a restatement of elementary economics: when something becomes more expensive, less of it will be purchased.'' --- Justice Stevens, dissenting, 468 U.S.\ at 788 (\href{https://supreme.justia.com/cases/federal/us/468/737/}{Justia}).
\end{quote}

\end{document}
